% Shared Manual / Manual II D.E.C. ledger source.
% Manual I branch is word-for-word the prior manual/sections/00_dec_ledger.tex content.
% Manual II branch renders the fractal-coordinate D.E.C. / SADAR / tunnelling ledger.
\ifdefined\AODManualII
\section{Pen-and-Paper AFC Raw Execution and D.E.C. Ledger for Fusion Scales}
\label{sec:pen-paper-dec-ledger}

AFC runs raw. The trace records the run. A$\Omega$D names detected motifs after the trace has been written. For Manual II this means that a fusion-scale fixture first declares the finite scope $B$, the region, the boundary, the admitted edge rows, and the update rule. The raw AFC run then emits D.E.C. rows. A$\Omega$D detectors read the trace and name retained motifs, formula closures, route-support events, contact/reclosure packets, SADAR contexts, and residual rows.

\aodnoteblock{D.E.C. ledger}{D.E.C. is the manual ledger format for finite AFC edge execution. It records the declared edge computation before target data or report maps are attached.}

\subsection{Fractal-coordinate D.E.C. row}

A Manual II D.E.C. row is read at a declared fractal coordinate and cut-frame:
\[
a\vdash B|t=(\partial F,\{t\}).
\]
The local tesseract witness is
\[
\epsilon\in Q_4=\{0,1\}^4,
\qquad
\mu\in\{-1,0,+1\}.
\]
The local cell is
\[
c=(a,\epsilon,\mu).
\]
A raw row is written
\[
\mathrm{D.E.C.}_e=(t,B,a,\epsilon,\mu,e,v_e,\sigma_e,\operatorname{adm}_e(B),w_e(B),P_e(B),\operatorname{route}_e(B)).
\]
The expanded CSV-facing schema for fractal-coordinate transparency is:
\begin{center}
\scriptsize
\begin{tabularx}{0.98\linewidth}{@{}p{0.24\linewidth}X@{}}
\toprule
field & role \\
\midrule
\texttt{tick} & cut-frame row index \\
\texttt{fractal\_coord} & retained address $a$ \\
\texttt{B} & declared boundary/window row \\
\texttt{epsilon\_Q4} & local $Q_4$ tesseract vertex $\epsilon$ \\
\texttt{mu} & ternary hinge fibre $\mu$ \\
\texttt{v}, \texttt{e}, \texttt{v\_e} & source, edge/event, and target node \\
\texttt{sigma\_e} & edge state \\
\texttt{adm\_e\_B} & admissibility on declared scope \\
\texttt{w\_e\_B}, \texttt{Z\_B} & declared integer weight and normalizer \\
\texttt{P\_num}, \texttt{P\_den}, \texttt{P\_exact} & exact normalized kernel fraction \\
\texttt{route\_e\_B} & route status or unit \\
\texttt{detector\_label} & A$\Omega$D trace label, when read \\
\texttt{adar\_id}, \texttt{sadar\_context\_id} & downstream ADAR/SADAR row identifiers, when promoted \\
\bottomrule
\end{tabularx}
\end{center}

The molecular-chain source ledger uses the compact columns
\[
(t,B,v,e,v_e,\sigma_e,\operatorname{adm}_e(B),w_e(B),P_e(B),\operatorname{route}_e(B)),
\]
and the demonstration ledgers expose the fractal-coordinate fields used for hand audit.

\subsection{$Q_4$ isotropic and anisotropic kernels}

For a local cell $c=(a,\epsilon,\mu)$, declare the admissible successor pool
\[
\mathcal A(c;B|t)=\{e:\operatorname{adm}_e(c;B|t)=1\}.
\]
The isotropic $Q_4$ kernel is
\[
P^{\rm iso}(e\mid c,B|t)=\frac{1}{|\mathcal A(c;B|t)|}.
\]
For a full unblocked $Q_4$ node,
\[
|\mathcal A|=4,
\qquad
P^{\rm iso}=(1/4,1/4,1/4,1/4).
\]
For an anisotropic kernel, declare integer weights before normalization:
\[
Z(c;B|t)=\sum_{e\in\mathcal A(c;B|t)}w_e(c;B|t),
\qquad
P^{\rm ani}(e\mid c,B|t)=\frac{w_e(c;B|t)}{Z(c;B|t)}.
\]
With
\[
w=(1,1,4,2),
\qquad
Z=8,
\]
one obtains
\[
P^{\rm ani}=(1/8,1/8,1/2,1/4).
\]

\subsection{Blocked hinge and hinge slide}

A blocked direct hinge has zero admissibility on the declared scope:
\[
\operatorname{adm}(e_0;B|t)=0.
\]
The slide pool is declared before normalization:
\[
\operatorname{Slide}_{\mu}(B|t)=\{(e_1,-1),(e_2,-1),(e_3,+1)\}.
\]
Declare post-block slide weights
\[
w'=(3,3,1),
\qquad
Z'=7.
\]
Then
\[
P_{\rm slide}=(3/7,3/7,1/7),
\]
with
\[
P_{\rm return}=3/7+3/7=6/7,
\qquad
P_{\rm out}=1/7.
\]

\subsection{Field tunnelling: hinge-slide window clip}
\label{subsec:field-tunnelling-window-clip}

Field tunnelling is the retained passage of a field route through a blocked hinge by slide-compatible successor support inside the declared window. Its compact spine is
\[
\begin{array}{c}
\text{blocked direct hinge}\to\text{declared slide pool}\to\text{exact slide kernel}\\[0.25em]
\to\text{window clip}\to\text{tunnel contribution}\to\text{boundary audit}.
\end{array}
\]
For a slide route $s_i$,
\[
\rho^D_{\omega,s_i}(B)=\min(RD_{s_i}(B),|\omega|_{s_i}),
\qquad
p^D_{s_i}(B)=C_{3,s_i}\rho^D_{\omega,s_i}(B).
\]
The scalar tunnel contribution is
\[
T^D_{\rm tunnel}(B|t)=\sum_{s_i\in\operatorname{Slide}_{\mu}(B|t)}P_{\rm slide}(s_i;B|t)p^D_{s_i}(B).
\]
With orientation retained, the corresponding flux contribution is
\[
\Phi^D_{\rm tunnel}(B|t)=\sum_{s_i\in\operatorname{Slide}_{\mu}(B|t)}P_{\rm slide}(s_i;B|t)p^D_{s_i}(B)A_{s_i}(B).
\]
For the worked row, declare the tunnel branch
\[
s_{\rm tunnel}=(e_3,+1),
\qquad
P_{\rm tunnel}=1/7.
\]
Declare
\[
RD_{s_{\rm tunnel}}=5,
\qquad
|\omega|_{s_{\rm tunnel}}=2,
\qquad
C_{3,s}=3.
\]
Then
\[
\rho^D_{\omega,s_{\rm tunnel}}=\min(5,2)=2,
\qquad
p^D_{s_{\rm tunnel}}=3\cdot 2=6.
\]
Thus
\[
T^D_{\rm tunnel}=(1/7)\cdot 6=6/7.
\]
\begin{center}
\scriptsize
\begin{tabularx}{0.96\linewidth}{@{}p{0.36\linewidth}X@{}}
\toprule
row & value \\
\midrule
direct hinge admissibility & $0$ \\
slide pool & $(e_1,-1),(e_2,-1),(e_3,+1)$ \\
slide weights & $3,3,1$ \\
slide probabilities & $3/7,3/7,1/7$ \\
tunnel branch & $(e_3,+1)$ \\
$RD_s$ & $5$ \\
$|\omega|_s$ & $2$ \\
$\rho^D_{\omega,s}$ & $2$ \\
$p^D_s$ & $6$ \\
$T^D_{\rm tunnel}$ & $6/7$ \\
\bottomrule
\end{tabularx}
\end{center}
The tunnel contribution is the window-clipped pressure carried by the slide-compatible outgoing successor.

\subsection{D.E.C. row-pair to ADAR and SADAR}

A routed D.E.C. row-pair supplies the declared returned-current pair. That pair is promoted to ADAR:
\[
ADAR_e(B)=(A_e(B),D_e(B),R^{\rm asym}_e(B)).
\]
The reflection-duration coupling is
\[
RCD_e(B)=D_e(B)\leftrightarrow R^{\rm asym}_e(B).
\]
The scalar window participation is
\[
\rho^D_{\omega,e}(B)=\min(\operatorname{dur}(RCD_e(B)),|\omega|_e),
\qquad
p^D_e(B)=C_{3,e}\rho^D_{\omega,e}(B).
\]
The SADAR flux contribution is
\[
\phi_e(B)=p^D_e(B)A_e(B),
\qquad
A_F(B)=\sum_{e\in F(B)}\phi_e(B).
\]
For a kernel-weighted expected readout over the D.E.C. admissible pool,
\[
A_F^{\rm exp}(B|t)=\sum_{e\in\mathcal A(c;B|t)}P_e(B|t)p^D_e(B)A_e(B).
\]
Thus $P_e$ is the AFC/D.E.C. kernel weight, $p^D_eA_e$ is the ADAR/SADAR content, and $P_ep^D_eA_e$ is the kernel-weighted SADAR expectation.

A compact hand example declares
\[
A_{12}=1/3,
\qquad
\rho^D_{\omega}=4,
\qquad
C_3=3.
\]
Then
\[
p^D=C_3\rho^D_{\omega}=12,
\qquad
S_B=p^D A_{12}=12(1/3)=4.
\]
Reverse orientation gives
\[
A_{21}=-1/3,
\qquad
S_B=-4.
\]

\subsection{Boundary closure and sheddic route audit}

Declare
\[
P^D=27,
\qquad
C_{\rm close}=24.
\]
Then
\[
\Delta_{\rm close}=P^D-C_{\rm close}=3,
\qquad
X_{\rm shedding}=\max(0,\Delta_{\rm close})=3.
\]
This is the positive surplus route audit used in the Dimonnanyro row:
\[
1{:}2{:}9_{\rm yro}=\text{Dimonnanyro}.
\]

\subsection{GAS tripeptide pen-and-paper raw row}

The chain \texttt{chain\_GAS\_tripeptide\_seed} has two admitted water-route support rows. The raw D.E.C. entries are:
\begin{center}
\scriptsize
\begin{tabularx}{0.98\linewidth}{@{}llllllllll@{}}
\toprule
\texttt{tick} & \texttt{B} & \texttt{v} & \texttt{e} & \texttt{v\_e} & \texttt{sigma\_e} & \texttt{adm\_e\_B} & \texttt{w\_e\_B} & \texttt{P\_e\_B} & \texttt{route\_e\_B} \\
\midrule
6 & \texttt{B\_mol\_005} & \texttt{v0} & \texttt{e0\_1} & \texttt{v1} & branch & 1 & 1 & $1/2$ & water\_equivalent\_route\_unit \\
7 & \texttt{B\_mol\_005} & \texttt{v1} & \texttt{e1\_2} & \texttt{v2} & branch & 1 & 1 & $1/2$ & water\_equivalent\_route\_unit \\
\bottomrule
\end{tabularx}
\end{center}
Thus the route probability and support count are computed by hand as
\[
P_{\mathrm{route}}=\frac12+\frac12=1,
\qquad
\mathrm{support\_units}=2.
\]

\subsection{Trace to detector to freeze}

The row chain is visible in the ledgers. Read it top to bottom:
\begin{center}
\scriptsize
\begin{tabular}{@{}c@{}}
\texttt{chain\_word\_spec.csv} \\
$\downarrow$ \\
\texttt{raw\_dec\_trace\_molecular\_chain.csv} \\
$\downarrow$ \\
\texttt{read\_only\_trace\_molecular\_chain.csv} \\
$\downarrow$ \\
\texttt{detected\_chain\_motifs.csv} \\
$\downarrow$ \\
\texttt{sadar\_detector\_context.csv} \\
$\downarrow$ \\
\texttt{chain\_formula\_predictions.csv} \\
$\downarrow$ \\
\texttt{molecular\_chain\_delta3\_audit.csv} \\
$\downarrow$ \\
\texttt{aod\_contact\_prediction\_freeze.csv} \\
$\downarrow$ \\
\texttt{protein\_contact\_score.csv}
\end{tabular}
\end{center}
For the GAS row this means:
\begin{align*}
\texttt{trace\_mol\_005} &\to \texttt{water\_route\_support\_accumulated} \to 2,\\
\texttt{motif\_mol\_005} &\to \texttt{water\_route\_chain\_closure},\\
\texttt{sadar\_mol\_005} &\to \texttt{water\_route\_reclosure\_context}.
\end{align*}
The freeze row is then computed as
\[
\mathrm{glycine}+\mathrm{alanine}+\mathrm{serine}-2H_2O
=\mathrm{C}_8\mathrm{H}_{15}\mathrm{N}_3\mathrm{O}_5.
\]
In named CHNOPS coordinates the check residual is
\[
\Delta_{\mathrm{CHNOPS}}=(0,0,0,0,0,0),
\qquad
\delta_3=0.
\]
Only after this freeze may the protein contact/reclosure packet cite
\[
\texttt{ChainWordSpec+ReadOnlyTrace+DetectedChainMotif+SADARContext}
\]
as its input basis.

\subsection{Manual II rendering status}

\aodnoteblock{Rendering status}{Manual II v40.02r08.6 is the compact worked-row rendering of the fusion-scale manual. The broader octave and target-lane source sections remain in the source tree and are tracked by the roadmap and data manifests.}

\aodnoteblock{Folding status}{This lane records sequence, contact-target, and contact/reclosure fixture rows. The coordinate-score registry carries RMSD, TM-score, GDT, and coordinate-level AOD prediction after a coordinate-level prediction freeze is declared.}

\aodnoteblock{Current score scope}{The current score rows are declared contact-residual fixture rows and deferred coordinate-score registries.}

\aodnoteblock{Ledger pointer}{The hand-auditable paths are summarized in \path{manual-2/data/dec/}.}

\else

\section{Pen-and-Paper AFC Raw Execution and D.E.C. Ledger}\label{manual:dec-ledger}
\aodnoteblock{Scope}{This section gives a hand-calculation path from one declared AFC edge computation into the A\(\Omega\)D motif names used later in the manual. It is a manual ledger section: it records exact edge weights, kernel fractions, route status, and boundary values.}
\aodprovenance{main:sec:afc-stokes-seed; main:subsec:q4-four-edge-implementation; main:subsec:hinge-slide; main:subsec:sadar; manual \secref{manual:exact-arithmetic-presentation}; manual \secref{manual:cross-return-short-window}.}

AFC runs raw. The trace records the run. A\(\Omega\)D names detected motifs. A simulation fixture first declares the AFC scope \(B\), boundary/region data, admissible channels, and update rule. The AFC run produces trace rows. A\(\Omega\)D detectors read the trace and classify motifs such as cross-return, yro route, duration-clipping, sheddic route, terminal split, Tau-ring support, orbital retention, and collision-support trace status.

D.E.C. means Declared Edge Computation ledger. In this manual section, D.E.C. is the manual ledger format used to expose one declared AFC edge computation. Use the dotted form D.E.C.; the ledger is a manual value container for declared edge data.

\[
\begin{aligned}
\text{AFC raw node}
&\to \text{D.E.C. ledger}
\to \text{exact kernel}
\to \text{blocked hinge}
\to \text{hinge slide},\\
&\to \text{oriented ADAR component}
\to \text{SADAR value}
\to \text{boundary audit}.
\end{aligned}
\]
Fixture-comparator rows form exact \(\delta_3\) tallies.

\subsection{D.E.C. row format}\label{manual:dec-ledger:row-format}
Each D.E.C. row records
\[
B,
\quad
v,
\quad
e,
\quad
v_e,
\quad
\sigma(e),
\quad
\operatorname{adm}(e;B),
\quad
w(e;B),
\quad
P(e;B),
\quad
\operatorname{route}(e;B).
\]
Here \(B\) is the declared boundary/window, \(v\) is the active node, \(e\) is an admissible edge slot, \(v_e\) is the Hamming--1 neighbor reached by that slot, \(\sigma(e)\in\{-1,0,+1\}\) is the local ternary state, \(\operatorname{adm}(e;B)\) records admissibility on the declared scope, \(w(e;B)\) is the declared integer weight, \(P(e;B)\) is the exact normalized kernel fraction, and \(\operatorname{route}(e;B)\) records returned, outgoing, blocked, slide, hinge, or pending route status.

\subsection{One-node D.E.C. ledger}\label{manual:dec-ledger:one-node}
Declare one \(Q_4\) node and four Hamming--1 edge slots:
\[
v=(0,0,0,0),
\qquad
\{e_1,e_2,e_3,e_4\}.
\]
Declare ternary edge states \(\sigma(e_i)\in\{-1,0,+1\}\). The following ledger records one declared anisotropic kernel.
\begin{table}[H]
\centering
\scriptsize
\caption{One-node D.E.C. ledger.}
\begin{tabular}{l c c c r c c}
\toprule
edge & neighbor & $\sigma(e)$ & admissible & weight & probability & route\\
\midrule
$e_1$ & $(1,0,0,0)$ & $+1$ & 1 & 1 & $1/8$ & outgoing\\
$e_2$ & $(0,1,0,0)$ & $-1$ & 1 & 1 & $1/8$ & returned\\
$e_3$ & $(0,0,1,0)$ & $0$ & 1 & 4 & $1/2$ & hinge\\
$e_4$ & $(0,0,0,1)$ & $+1$ & 1 & 2 & $1/4$ & outgoing\\
\bottomrule
\end{tabular}
\label{tab:manual-dec-one-node-ledger}
\end{table}

\subsection{Exact kernel normalization}\label{manual:dec-ledger:kernel}
For isotropic admissible branching over four successors,
\[
P^{\mathrm{iso}}(e_i;B)=\frac14.
\]
For the declared anisotropic weights,
\[
w=(1,1,4,2),
\qquad
Z=1+1+4+2=8.
\]
The anisotropic kernel is
\[
P^{\mathrm{ani}}=(1/8,1/8,1/2,1/4).
\]
\begin{table}[H]
\centering
\scriptsize
\caption{Exact anisotropic kernel fractions.}
\begin{tabular}{l r c}
\toprule
edge & weight & probability\\
\midrule
$e_1$ & 1 & $1/8$\\
$e_2$ & 1 & $1/8$\\
$e_3$ & 4 & $1/2$\\
$e_4$ & 2 & $1/4$\\
\bottomrule
\end{tabular}
\label{tab:manual-dec-kernel-fractions}
\end{table}
\aodliteraturenote{Markov kernels and random walks: Norris~\cite{norrisMarkov}, Levin--Peres--Wilmer~\cite{levinPeresWilmer}, Lawler--Limic~\cite{lawlerLimic}.}

\subsection{Blocking rule}\label{manual:dec-ledger:blocking-rule}
Blocking is zero admissibility. Once same-resolution hinge dwell is blocked on the declared scope, the hinge-state row is removed from the successor pool and the kernel is renormalized over branch-oriented continuations.
\begin{table}[H]
\centering
\scriptsize
\caption{D.E.C. blocking and slide rules.}
\begin{tabular}{l l}
\toprule
rule & meaning\\
\midrule
isotropic & equal probability over admissible successors\\
anisotropic & unequal declared weights over admissible successors\\
blocked & zero admissibility on the declared scope\\
hinge slide & renormalization after same-resolution dwell is removed\\
wave trace & repeated node updates after hinge slide\\
\bottomrule
\end{tabular}
\label{tab:manual-dec-blocking-rules}
\end{table}

\subsection{Blocked hinge and hinge slide}\label{manual:dec-ledger:hinge-slide}
In the following block event, the prior hinge-state row \((e_3,0)\) attempts same-resolution dwell, but the declared scope sets \(\operatorname{adm}(e_3,0;B)=0\). The hinge-state row \((e_3,0)\), not the edge slot \(e_3\) itself, is removed before renormalization; the \(e_3\) slot may re-enter only as an admitted branch state. Declare the post-block slide pool
\[
\{(e_1,-1),(e_2,-1),(e_3,+1)\}
\]
with weights
\[
w'((e_1,-1),(e_2,-1),(e_3,+1))=(3,3,1),
\qquad
Z'=3+3+1=7.
\]
Then
\[
P(e_1,-1)=3/7,
\qquad
P(e_2,-1)=3/7,
\qquad
P(e_3,+1)=1/7.
\]
The returned and outgoing slide fractions are
\[
P_{\mathrm{returned}}=3/7+3/7=6/7,
\qquad
P_{\mathrm{outgoing}}=1/7.
\]


\subsection{Field tunnelling: hinge-slide window clip}\label{manual:dec-ledger:field-tunnelling}
\aodnoteblock{Data-support class}{Field tunnelling is a D0 exact internal fixture. It uses declared hinge-slide admissibility, exact slide-kernel fractions, window-clipped duration participation, and scalar tunnel contribution.}

Field tunnelling is the retained passage of a field route through a blocked hinge by slide-compatible successor support inside the declared window. The manual spine is
\[
\begin{aligned}
\text{blocked direct hinge}
&\to \text{declared slide pool}
\to \text{exact slide kernel}
\to \text{window clip},\\
&\to \text{tunnel contribution}
\to \text{boundary audit}.
\end{aligned}
\]
Fix the declared boundary/window
\[
B=(\partial F,\omega).
\]
At cut-frame \(B|t\), declare a direct hinge route \(e_0\) with
\[
\operatorname{adm}(e_0;B|t)=0.
\]
Declare the slide pool
\[
\operatorname{Slide}_{\mu}(B|t)=\{s_1,\ldots,s_n\}.
\]
Normalize only after the slide pool is declared:
\[
Z_{\mathrm{slide}}(B|t)
=
\sum_{s_i\in\operatorname{Slide}_{\mu}(B|t)}w(s_i;B|t),
\qquad
P_{\mathrm{slide}}(s_i;B|t)
=
\frac{w(s_i;B|t)}{Z_{\mathrm{slide}}(B|t)}.
\]
For each slide-compatible successor,
\[
\rho^D_{\omega,s_i}(B)=\min(RD_{s_i}(B),|\omega|_{s_i}),
\qquad
p^D_{s_i}(B)=C_{3,s_i}\rho^D_{\omega,s_i}(B).
\]
The scalar tunnel contribution is
\[
T^D_{\mathrm{tunnel}}(B|t)
=
\sum_{s_i\in\operatorname{Slide}_{\mu}(B|t)}
P_{\mathrm{slide}}(s_i;B|t)\,p^D_{s_i}(B).
\]
When oriented attention is declared, the oriented tunnel flux is
\[
\Phi^D_{\mathrm{tunnel}}(B|t)
=
\sum_{s_i\in\operatorname{Slide}_{\mu}(B|t)}
P_{\mathrm{slide}}(s_i;B|t)\,p^D_{s_i}(B)\,A_{s_i}(B).
\]
Using the hinge-slide pool from \secref{manual:dec-ledger:hinge-slide}, declare the tunnel branch
\[
s_{\mathrm{tunnel}}=(e_3,+1),
\qquad
P_{\mathrm{tunnel}}=1/7.
\]
Declare
\[
RD_{s_{\mathrm{tunnel}}}=5,
\qquad
|\omega|_{s_{\mathrm{tunnel}}}=2,
\qquad
C_{3,s}=3.
\]
Then
\[
\rho^D_{\omega,s_{\mathrm{tunnel}}}=\min(5,2)=2,
\qquad
p^D_{s_{\mathrm{tunnel}}}=3\cdot2=6,
\]
and
\[
T^D_{\mathrm{tunnel}}=(1/7)\cdot6=6/7.
\]
The tunnel contribution is the window-clipped pressure carried by the slide-compatible outgoing successor.
\begin{table}[H]
\centering
\scriptsize
\caption{Field tunnelling hinge-slide window-clip fixture.}
\begin{tabular}{l l}
\toprule
row & value\\
\midrule
direct hinge admissibility & $\operatorname{adm}(e_0;B|t)=0$\\
slide pool & $\{(e_1,-1),(e_2,-1),(e_3,+1)\}$\\
slide weights & $(3,3,1)$\\
slide probabilities & $(3/7,3/7,1/7)$\\
tunnel branch & $(e_3,+1)$\\
$P_{\mathrm{tunnel}}$ & $1/7$\\
$RD_s$ & $5$\\
$|\omega|_s$ & $2$\\
$\rho^D_{\omega,s}$ & $2$\\
$C_{3,s}$ & $3$\\
$p^D_s$ & $6$\\
$T^D_{\mathrm{tunnel}}$ & $6/7$\\
\bottomrule
\end{tabular}
\label{tab:manual-field-tunnelling-fixture}
\end{table}

\subsection{D.E.C. row-pair to ADAR and SADAR}\label{manual:dec-ledger:adar-sadar}
A routed D.E.C. row-pair supplies the declared returned-current pair. That pair is promoted to an oriented ADAR component on the declared boundary/window. Declare
\[
A_{12}=1/3,
\qquad
\rho^D_\omega=4,
\qquad
C_3=3.
\]
Then
\[
p^D=C_3\rho^D_\omega=12,
\qquad
\SADARop_B=p^D A_{12}=12(1/3)=4.
\]
Reverse orientation gives
\[
A_{21}=-1/3,
\qquad
\SADARop_B=-4.
\]

\subsection{Boundary closure and sheddic route audit}\label{manual:dec-ledger:boundary-audit}
Declare
\[
P^D=27,
\qquad
C_{\mathrm{close}}=24.
\]
Then
\[
\Delta_{\mathrm{close}}=P^D-C_{\mathrm{close}}=3,
\qquad
X_{\mathrm{shedding}}=\max(0,\Delta_{\mathrm{close}})=3.
\]
This is the positive surplus route audit used in the Dimonnanyro row:
\[
1{:}2{:}9_{\mathrm{yro}}=\mathrm{Dimonnanyro}.
\]
Fixture-comparator rows record exact \(\delta_3\) residuals. Boundary-route rows record boundary values and route audits.


\subsection{D.E.C. audit invariants}\label{manual:dec-ledger:audit-hooks}
A D.E.C. row records
\[
B,
\quad
v,
\quad
e,
\quad
v_e,
\quad
\sigma(e),
\quad
\operatorname{adm}(e;B),
\quad
w(e;B),
\quad
P(e;B),
\quad
\operatorname{route}(e;B).
\]
The normalized kernel satisfies
\[
\sum_e P(e;B)=1.
\]
Zero admissibility gives zero probability:
\[
\operatorname{adm}(e;B)=0\Rightarrow P(e;B)=0.
\]
A hinge-slide row records
\[
\mathrm{Adm}^0_B(e,\sigma)=\varnothing
\]
and the renormalized branch-continuation kernel. Blocking is zero admissibility.

A SADAR row records an oriented ADAR pair and satisfies
\[
\SADARop_B(C_{\bar e})=-\SADARop_B(C_e).
\]

Fixture-comparator rows form exact \(\delta_3\) tallies. A Tau missing-burden row carries integer provenance
\[
\mathrm{miss}_\tau(K)=M_{\mathrm{miss}}(K)/M_{\mathrm{ref}}(K),
\]
with any decimal shown only as a display coordinate.

\subsection{Tau missing-burden exact ratio}\label{manual:dec-ledger:tau-ratio}
The Tau missing-burden diagnostic is recorded with exact provenance:
\[
\mathrm{miss}_\tau(K)=\frac{M_{\mathrm{miss}}(K)}{M_{\mathrm{ref}}(K)}.
\]
For the displayed Dimonnanyro row,
\[
M_{\mathrm{miss}}=31,
\qquad
M_{\mathrm{ref}}=50,
\qquad
\mathrm{miss}_\tau(1{:}2{:}9_{\mathrm{yro}})=\frac{31}{50}.
\]
The decimal presentation field is
\[
\mathrm{missing\_fraction\_display}=0.62.
\]
\begin{table}[H]
\centering
\scriptsize
\caption{Tau missing-burden exact-ratio fields.}
\begin{tabular}{l l}
\toprule
field & meaning\\
\midrule
missing\_num & numerator $M_{\mathrm{miss}}$\\
missing\_den & denominator $M_{\mathrm{ref}}$\\
missing\_frac\_exact & exact ratio\\
missing\_fraction\_display & decimal presentation field\\
\bottomrule
\end{tabular}
\label{tab:manual-dec-tau-missing-ratio-fields}
\end{table}

\fi
